% ---------------------------------------------------------------------
% Manuscript Sec. VI-E, revised around the envelope-witness chain.
%
% Numbering map against the current manuscript:
%   (90) envelope, (91) rho_bar, (92) sampled record + LS objective
%        -- kept as they stand; (92) is the only place the sample
%        index i remains.
%   (93)          -- NEW: the witness chain (replaces the old
%                    intermediate steps).
%   (94)          -- NEW: the window RMS of the envelope, B(x).
%   (95)-(97c)    -- kept, with the index i dropped (function-level
%                    identities); they feed VI-F unchanged.
%   (98)-(108)    -- DELETED (the delta-e cap and Delta_pre; see the
%                    closing remark).
%   (98') = old (109) noise model, renumbered; (99') the headline
%   bound, replacing old (110).  Downstream equations shift by -9.
%
% The text below is manuscript-ready two-column prose; this preview
% wraps it in a standalone article for compilation.  Figure labels
% \ref{fig:bound_sim}, \ref{fig:bound_hw} refer to
% fig_bound_final_sim.pdf and fig_bound_final_hw.pdf.
% ---------------------------------------------------------------------

\documentclass[10pt]{article}
\usepackage[a4paper,margin=22mm,twocolumn]{geometry}
\usepackage{amsmath,amssymb,booktabs}
\setlength{\parskip}{0pt}\setlength{\parindent}{1em}
\newcommand{\Mdot}{\dot M}
\newcommand{\rbar}{\bar\rho}
\newcommand{\rms}{\mathrm{RMS}}
\newcommand{\Wz}{W\!z_{CoM}}
% preview stand-ins for the manuscript figure labels; in the
% manuscript use \ref{fig:bound_sim} / \ref{fig:bound_hw} on the two
% validation figures (fig_bound_final_sim.pdf, fig_bound_final_hw.pdf)
\newcommand{\figboundsim}{15}
\newcommand{\figboundhw}{16}
\begin{document}
\twocolumn[{\centering\Large\bfseries
Sec.~VI-E, revised: RMS residual bound of the PNLS fit\\[2mm]
\normalsize\normalfont preview build; see the header comment for the
numbering map\\[6mm]}]

\subsection*{E. RMS residual bound of the PNLS fit}

The envelope (90) bounds the deviation $e_\omega$ pointwise; this
subsection converts it into a certificate on the quantity the
pipeline actually reports, the residual RMS of the deployed fit. The
argument needs one observation and one integral.

\paragraph*{The witness inequality}
The record in the window is
$y(\tau) = \omega_{\rm nom}(\tau) + e_\omega(\tau) + b + n(\tau)$,
with $\tau$ counted from the \emph{realised} onset $t_{\rm on}$ ---
the instant the restoring moment reverses at the vehicle, wherever
the actuation chain places it relative to the commanded ramp clock.
The PNLS stage fits the family
$\hat h(\tau;\theta) = C_1^*\big(\cosh(C_2^*(\tau - t^*)) -
1\big)\mathbf 1[\tau \ge t^*] + C$ with
$\theta = (C_1^*, C_2^*, t^*)$ free and the baseline $C$ pinned to
the pre-onset segment median --- the convention of Algorithm~1, so
the two stages share it --- minimising the objective of (92). The
nominal solution is itself a member of this family: the choice
$\theta_0 = (C_1, C_2, t_{\rm on})$ reproduces
$\omega_{\rm nom} + C$ with $C$ estimating the bias $b$ to
$O(\sigma_n/\sqrt{N_{\rm pre}})$, a discrepancy that rides in the
noise term. A minimum never exceeds the value at one feasible
point, so the fit cannot lose to this \emph{witness}, and
\begin{equation}
  \rms(r) \;\le\; \rms(e_\omega + n)
  \;\le\; \rms(e_\omega) + \rms(n).
  \tag{93}
\end{equation}
No property of the minimiser is used --- the fitted parameters may
wander along the amplitude--onset ridge of Sec.~V --- and no
admissibility condition arises, because the witness is the nominal
itself rather than a shifted surrogate. Two consequences are worth
stating. First, a pure actuation delay shifts $t_{\rm on}$ and
nothing else; since $t^*$ is free, the witness adopts the realised
onset and the delay never enters the bound --- no latency term is
required. Second, before the onset both the witness and the record
sit at the baseline, so the pre-onset segment contributes only
noise: no pre-onset penalty of the form (108) appears, and
extending the RMS to the full window only dilutes the left side of
(93).

\paragraph*{The window RMS of the envelope}
Squaring the envelope (90) and integrating over the window, with
$x = C_2\tau_{\rm end}$,
\begin{equation}
  \rms(e_\omega) \;\le\;
  \rbar K C_2\sqrt{\frac{B(x)}{x}},
  \qquad
  B(x) = \frac{\sinh 2x}{4} - \frac{x}{2},
  \tag{94}
\end{equation}
$B(x) = \int_0^x \sinh^2 u\,du$ by the double-angle identity;
$\sqrt{B(x)/x}$ is the window RMS of the growing mode
$\sinh(C_2\tau)$, and $x$ follows per ramp rate from the design
sweep $\sinh x - x = \varphi_{\max}\Wz C_2/\Mdot$ of Sec.~VI-D. The
tilt cap $\varphi_{\max}$ must dominate the tilt the windows
actually reach, and is set per campaign: $5^\circ$ in simulation,
where the gate stops the window near $4^\circ$, and $8^\circ$ on
hardware, where the slightly uneven ground of the test site lets
the in-window excursion reach $7.7^\circ$. Every appearance of
$\varphi_{\max}$ is monotone, so a dominating cap only over-covers.

\paragraph*{The measured decomposition and the onset shift}
For the critical-moment readout of Sec.~VI-F the deviation is split
against the growing mode,
\begin{equation}
  e_\omega(\tau) \;=\; \beta\sinh(C_2\tau) + \delta e_\omega(\tau),
  \tag{95}
\end{equation}
with the end-matched coefficient
\begin{equation}
  \beta \;=\;
  \frac{e_\omega(\tau_{\rm end})}{\sinh(C_2\tau_{\rm end})},
  \tag{96}
\end{equation}
so that $\delta e_\omega$ vanishes at both window ends. [(95)--(97)
are function-level identities in $\tau$; the sample index $i$
belongs to (92) alone.] The envelope (90) caps the coefficient a
priori,
\begin{equation}
  |\beta| \;\le\; \frac{\rbar}{J_PC_2},
  \tag{97a}
\end{equation}
and the $\sinh$ multiple is exactly the tangent direction of the
onset shift, so the fitted branch absorbs it as a critical-time
offset
\begin{equation}
  \tanh(C_2\,\delta t_c) \;=\; -\frac{\beta}{C_1},
  \qquad \delta t_c < 0,
  \tag{97b}
\end{equation}
which Sec.~VI-F converts into the critical-moment error
$|\Delta M_{\rm crit}| = \Mdot|\delta t_c|$. This absorption is also
why the measured residuals stay flat as the ramp rate grows: the
dominant, exponentially growing part of $e_\omega$ is precisely the
component the family's onset freedom removes.

\paragraph*{The noise term}
What the family cannot represent and the dynamics do not produce
remains in the residual as noise. On hardware the rotor-induced
vibration has no pre-experiment model, so its amplitude is read
from the run's own out-of-band content: with $n_{\rm hi}$ the
residual RMS above $f_c = 5$ Hz --- a band the cosh family cannot
occupy, its fastest scale being $C_2/2\pi \approx 0.8$ Hz --- the
full-band estimate is
\begin{equation}
  \rms(n) \;=\; \rms(n_{\rm hi})
  \sqrt{1 + \big(s_{\rm med}\,\kappa_q\big)^2},
  \tag{98}
\end{equation}
with the quiet-segment band ratio and in-window enhancement of
Sec.~VI-D ($s_{\rm med}\kappa_q = 1.31$ on the campaign). In
simulation the sensor noise is declared, Gaussian with
$\sigma_n = 2.450\times10^{-4}$ rad/s $= 0.014^\circ$/s --- an
order below the smallest envelope --- so the noise term is dropped
there.

\paragraph*{The bound, and its validation}
Combining (93), (94) and (98),
\begin{equation}
  \rms(r) \;\le\;
  \rbar K C_2\sqrt{\frac{B(x)}{x}} \;+\; \rms(n),
  \tag{99}
\end{equation}
with every quantity on the right a pre-experiment design constant
except the hardware $n_{\rm hi}$. Figs.~\figboundsim{} and
\figboundhw{} test (99) on every run of both campaigns, the
measured side being the deployed free fit exactly as it runs in the
pipeline --- both shape constants and the onset free, full window,
no calibration on either side. The cap is evaluated once per ramp
rate at the worst corner of the design box --- both leg arms, the
far corner of the offset rectangle, both campaign masses --- so the
theory enters each figure as a single curve, with no per-run
knowledge on the theory side. The bound is a conditional guarantee
--- offset inside the design rectangle, tilt under
$\varphi_{\max}$, channels as modelled --- and on every run
satisfying its hypothesis it holds without exception: in simulation
($\varphi_{\max} = 5^\circ$, envelope alone) $252/252$ in-box runs
sit under the curve (worst used fraction $0.79$), the two
exceedances confined to S9 and S11, the deliberate out-of-box
probes whose offsets $(32,32)$ and $(38,14)$ mm lie beyond the
rectangle --- the certificate fails exactly where its own premise
does. On hardware ($\varphi_{\max} = 8^\circ$, envelope plus the
campaign's largest vibration term $\max n_{\rm hi}\sqrt{1 +
(s_{\rm med}\kappa_q)^2} = 2.47^\circ$/s in (98)) all $140$ runs
sit under the curve, the worst using $0.55$ of its cap. The median used fraction
is $0.20$--$0.28$ on both campaigns: one construction, validated on
both platforms, the panels differing exactly in the channels the
platforms differ in.

\paragraph*{Remark}
An earlier version of this analysis bounded $\delta e_\omega$
itself through the Dirichlet response of the forcing channels
[the former (98)--(108)]. That cap is exact on the dynamics it
models --- on the exactly integrated tip-over it covers
$\delta e_\omega$ at every ramp rate with $2.5$--$3\times$ margin
--- but $\delta e_\omega$ is defined against the true nominal and
therefore cannot be measured separately on a platform: any error
in the pinned constants or onset injects in-family content into
the measured difference, which dominates it in practice. A
platform-validated certificate must not depend on knowing the true
constants on the measured side, which is exactly what (93)--(99)
provide; the Dirichlet machinery is therefore omitted from the
validation. The decomposition (95)--(97b) retains its two roles:
the a-priori ceiling on the onset shift consumed by Sec.~VI-F, and
the explanation of the residual's rate-flatness.

\end{document}
